%%% SECTION 7: PHYSICAL AI STANDARD LEVEL AND UNIFICATION STANDARD LEVEL %%%

[This section should be approximately 10-12 pages. The major theme is that the
PSL (Physical AI Standard Level) and USL (Unification Standard Level) frameworks
complement each other to create a comprehensive, quantitative standards system
for responsible robot usage in oncology clinical trials. PSL provides the
regulatory compliance dimensions for trial site operations, while USL evaluates
individual robot readiness. Together, they ensure that both the site
infrastructure and the robots operating within it meet the standards required
for safe, effective Physical AI oncology trials.]

[PRIMARY SOURCES: national-platform/usl\_standard/ directory containing
01\_preamble\_introduction\_framework.tex and 02\_results\_discussion\_closing.tex
for the USL framework; and national-platform/new\_trial\_psl/ directory
containing all 11 files (01\_preamble\_and\_sb1042.tex through
11\_emergency\_preparedness.tex) for the PSL framework embedded within the
clinical trial site documentation.]

\subsection{Overview: Complementary Standards}
\label{subsec:standards-overview}

[Explain how PSL and USL serve different but complementary purposes:
PSL evaluates site-level regulatory compliance across 3 dimensions, while
USL evaluates individual robot readiness across 4 dimensions. Together, a
trial site can only operate when both PSL requirements are met (site is
compliant) and USL scores are sufficient (robots are ready). This dual-layer
approach is unique to this National Platform and more thorough than any
existing FDA guidance on robot or AI evaluation.]

[MAJOR THEME: PSL and USL together enable a quantitative, reproducible
assessment of Physical AI trial readiness that replaces subjective evaluation
with objective scoring. This is essential for scaling Physical AI oncology
trials nationwide, as each new site can be evaluated against the same
standards.]

\subsection{Physical AI Standard Level (PSL) Framework}
\label{subsec:psl-framework}

[Build from national-platform/new\_trial\_psl/ files, particularly
01\_preamble\_and\_sb1042.tex which establishes the PSL framework context.
The PSL framework was developed in conjunction with the 24-hour clinical trial
simulation evidence.]

\subsubsection{The Three PSL Dimensions}
[Detail the three regulatory dimensions of PSL and explain what each
dimension measures. Draw from the PSL framework as embedded in the clinical
trial site documentation. Each dimension should be defined with its scope,
measurement criteria, and passing thresholds.]

[MAJOR THEME: The three PSL dimensions are important for responsible robot
usage because they ensure that a trial site's regulatory environment,
operational procedures, and safety infrastructure all meet minimum standards
before any Physical AI procedures are performed. This prevents scenarios where
technically capable robots operate in insufficiently prepared environments.]

\subsubsection{PSL Scoring Methodology}
[Explain how PSL scores are computed, including the weighting of each
dimension, the criteria for each score level, and the minimum overall score
required for trial site activation. Reference the 24-hour simulation
evidence from Clinical Trial Site Complete Documentation \cite{site-docs}
as validation of the PSL framework.]

\subsubsection{PSL and Site Activation}
[Explain how PSL scores determine whether a trial site can be activated for
Physical AI operations. Cover the relationship between PSL compliance and
the site activation procedures detailed in
national-platform/new\_trial\_psl/10\_activation\_sops.tex.]

\subsection{Unification Standard Level (USL) Framework}
\label{subsec:usl-framework}

[Build from national-platform/usl\_standard/01\_preamble\_introduction\_framework.tex,
covering the introduction, need for standardized robot readiness, and the
USL scoring methodology.]

\subsubsection{The Four USL Dimensions}
[Detail each of the four USL dimensions: regulatory compliance, clinical
integration, technical maturity, and safety architecture. For each dimension,
provide the definition, measurement criteria, scoring range, and significance
for Physical AI oncology trials.]

[MAJOR THEME: The four USL dimensions are important for responsible robot usage
because they evaluate each robot type holistically rather than on a single
criterion. A robot may be technically advanced but lack regulatory compliance
documentation, or may have excellent safety features but insufficient clinical
integration capabilities. USL ensures all four dimensions are evaluated.]

\subsubsection{USL Scoring Methodology}
[Build from national-platform/usl\_standard/01\_preamble\_introduction\_framework.tex,
Section 2 (Methods). Cover the development process, dimension definitions,
score computation formulas, bands and levels, category-specific scoring
engines, and individual robot evaluation modules.]

\subsubsection{USL Score Computation}
[Detail the mathematical framework for USL score computation, including how
individual dimension scores are combined into an overall USL score, the
band/level classification system, and the interpretation of results.]

\subsection{USL Results Across Robot Categories}
\label{subsec:usl-results}

[Build from national-platform/usl\_standard/02\_results\_discussion\_closing.tex,
Section 3 (Results). Present the complete USL scores table covering all ten
robot categories evaluated. Include:]

\subsubsection{Surgical Robot Results}
[Cover USL scores for surgical robots. Reference A Cancer Patient's Journey,
Physical AI \cite{patient-journey-paper} for clinical context of surgical
robot usage.]

\subsubsection{Cobot and Humanoid Results}
[Cover USL scores for cobots and humanoid robots. Reference Patient
Instructions: Physical AI Trials \cite{patient-instructions-paper} for patient
interaction context.]

\subsubsection{Specialized Robot Results}
[Cover USL scores for radiotherapy positioning robots, needle-placement robots,
social companion robots, imaging robots, steerable needle robots, and
rehabilitation exoskeletons. Reference the patient robot instructions from
national-platform/patient\_robot/ files for detailed descriptions of each
robot type.]

\subsubsection{Cross-Category Comparison}
[Present the cross-category comparison of USL scores, identifying which robot
types score highest on each dimension and overall. Discuss what the variation
in scores means for Physical AI trial readiness and where targeted improvement
is needed.]

\subsection{How PSL and USL Complement Each Other}
\label{subsec:psl-usl-complement}

[MAJOR THEME: Synthesize how the 3 PSL dimensions and 4 USL dimensions work
together. PSL ensures the environment is ready; USL ensures the robots are
ready. A Physical AI trial site requires both: a compliant, well-prepared
facility (PSL) populated with properly qualified robots (USL). Provide examples
of how a site might pass PSL but fail USL for certain robot types, or how a
robot might have high USL scores but be unable to operate at a site that fails
PSL. This complementary relationship makes the standards framework more robust
than either standard alone.]

[Reference Clinical Trial Site Complete Documentation \cite{site-docs} and
Physical AI Unification Standard Level \cite{usl-standard} as the primary
implementations of PSL and USL respectively. Connect to the MCP conformance
levels from Physical AI National MCP Servers \cite{national-mcp-paper} as an
additional interoperability standard that interfaces with both PSL and USL.]

\subsection{Impact on Nationwide Deployment}
\label{subsec:standards-impact}

[Build from national-platform/usl\_standard/02\_results\_discussion\_closing.tex,
Section 4 (Discussion). Cover the open-source ecosystem, hardware vs.
readiness distinctions, clinical trial readiness implications, and category
differences. Explain how standardized PSL and USL scores enable efficient
nationwide scaling of Physical AI oncology trials by providing a common
evaluation framework that all sites and robots must meet.]

[Reference the three code repositories \cite{main-repo, mcp-repo, fl-repo}
as open-source implementations that support PSL and USL assessment.]
